\uuid{UZD9}
\titre{Résolution d'équation différentielle du premier ordre}
\auteur{bodin ; croizat ; sacre}
\organisation{exo7}
\licence{CC BY-NC-SA 4.0}
\datecreate{2022-03-08}
\niveau{L1}
\module{Analyse}
\chapitre{Equation différentielle}
\sousChapitre{Résolution d'équation différentielle du premier ordre}
\difficulte{2}
\theme{}
\exo7qcmid{hilisit.equations-differentielles.030}

\contenu{
\begin{question}
\qclassification{225.01}
Que peut-on dire des solutions de l'équation différentielle $y'=ay$ ?
\begin{answers}
  \bad{Ce sont toutes des fonctions croissantes sur $\Rr$.}
  \bad{Ce sont toutes des fonctions décroissantes sur $\Rr$.}
  \bad{Si $a\ge 0$, ce sont des fonctions croissantes sur $\Rr$.}
  \good{Ce sont toutes des fonctions monotones sur $\Rr$.}
\end{answers}
\begin{explanations}
Les solutions de l'équation différentielle $y'=ay$ sont les fonctions $C e^{ax}$ avec $C$ constante réelle. Si $a\ge 0$, ce sont des fonctions croissantes pour $C\ge 0$ et décroissantes pour $C\le 0$. Si $a\le 0$, ce sont des fonctions décroissantes pour $C\ge 0$ et croissantes pour $C\le0$. Dans tous les cas, ce sont toutes des fonctions monotones sur $\Rr$.
\end{explanations}
\end{question}
}
